\documentclass[11pt]{article}
\usepackage[margin=1in]{geometry}
\usepackage{amsmath,amssymb,amsthm}
\usepackage{parskip}

\newtheorem{theorem}{Theorem}
\newtheorem{lemma}[theorem]{Lemma}
\newtheorem{corollary}[theorem]{Corollary}
\newtheorem{proposition}[theorem]{Proposition}
\theoremstyle{definition}
\newtheorem{definition}[theorem]{Definition}
\theoremstyle{remark}
\newtheorem{remark}[theorem]{Remark}

\title{Disproof of conjecture \texttt{00000001097}}
\author{earthking11}
\date{2026-09-12}

\begin{document}
\maketitle

\section*{Verdict: FALSE}

We disprove conjecture \texttt{00000001097} as stated. The refutation is
unconditional and elementary. The mean root count of $x^d + ax + b$ over
$\mathbb{F}_q$ is \emph{exactly} $1$ for every prime power $q$ and every $d$; the
claimed ``main term'' is strictly larger than $1$ and, when
$\gcd(d-1,q-1)=1$, coincides with the second moment $2-1/q$ rather than with the
mean; and the conjectured variance $2-1/q$ is off by exactly $1$. At $q=5$, $d=2$
the conjecture asserts both $9/5$ and $1$ for the same mean, so it is internally
inconsistent.

\section{The conjecture}

\begin{definition}[as stated]
The \emph{mean root count} of the trinomial $x^d + ax + b$ over $\mathbb{F}_q$ is
the expectation of the number of roots, taken over uniform
$(a,b) \in \mathbb{F}_q^2$.
\end{definition}

\begin{quote}
\textbf{Conjecture.} The mean equals $1 + (q-1)/q^{\gcd(d-1,q-1)}$ as the main
term; when $\gcd(d-1,q-1) = 1$ the mean is exactly $1$, and the variance of the
mean is $2 - 1/q$.
\end{quote}

Throughout, $q$ is a prime power and $\mathbb{F}_q$ the field with $q$ elements.
For $(a,b) \in \mathbb{F}_q^2$ put
\[
N(a,b) \;=\; \#\{\,x \in \mathbb{F}_q : x^d + ax + b = 0\,\},
\]
so that the mean root count is
\[
M(q,d) \;=\; \frac{1}{q^2} \sum_{a,b \in \mathbb{F}_q} N(a,b).
\]
We write $g = \gcd(d-1, q-1)$.

\section{The exact mean is $1$}

\begin{theorem}\label{thm:mean-one}
For every prime power $q$ and every integer $d \ge 1$,
\[
M(q,d) = 1.
\]
\end{theorem}

\begin{proof}
We count the triples $(a,b,x) \in \mathbb{F}_q^3$ satisfying
$x^d + ax + b = 0$. By definition of $M$,
\[
q^2 \, M(q,d) \;=\; \#\{(a,b,x) \in \mathbb{F}_q^3 : x^d + ax + b = 0\}.
\]
Fix $x \in \mathbb{F}_q$ and $a \in \mathbb{F}_q$. The map
\[
\mathbb{F}_q \to \mathbb{F}_q, \qquad b \mapsto x^d + ax + b
\]
is a bijection: it is translation by the constant $x^d + ax$, with inverse
$c \mapsto c - (x^d+ax)$. Hence it has exactly one zero, namely
$b = -x^d - ax$. Consequently, for each of the $q$ choices of $x$ and each of the
$q$ choices of $a$ there is exactly one $b$ for which the equation holds. The
number of triples is therefore $q \cdot q = q^2$, and
\[
M(q,d) \;=\; \frac{q^2}{q^2} \;=\; 1. \qedhere
\]
\end{proof}

Only the field property of $\mathbb{F}_q$ is used, via invertibility of translation
by $b$; no hypothesis on $d$ is needed.

\begin{corollary}\label{cor:main-term-wrong}
For every prime power $q>1$ and every $d$, the conjectured value
\[
1 + \frac{q-1}{q^{\gcd(d-1,q-1)}}
\]
is strictly greater than $1$, and therefore does not equal $M(q,d)$.
\end{corollary}

\begin{proof}
Since $q>1$ we have $q-1>0$, and $q^{g} \ge q > 0$ for $g = \gcd(d-1,q-1) \ge 1$.
Hence $(q-1)/q^{g} > 0$, so the displayed quantity is $>1$. By
Theorem~\ref{thm:mean-one}, $M(q,d)=1$, so equality fails.
\end{proof}

\section{Internal contradiction at $q=5$, $d=2$}

\begin{theorem}\label{thm:contradiction}
The conjecture is internally inconsistent.
\end{theorem}

\begin{proof}
Take $q = 5$ and $d = 2$. Then
\[
\gcd(d-1,\,q-1) = \gcd(1,4) = 1.
\]
The conjecture asserts, within one sentence, both
\[
M(5,2) = 1 + \frac{5-1}{5^{1}} = 1 + \frac{4}{5} = \frac{9}{5}
\qquad\text{and}\qquad
M(5,2) = 1 \quad\text{(``exactly $1$'' when $\gcd(d-1,q-1)=1$).}
\]
Since $9/5 \ne 1$, these two clauses cannot both hold. The same argument applies to
every $q>1$ with $\gcd(d-1,q-1)=1$: for instance $q=7$, $d=2$ gives
$\gcd(1,6)=1$ and $1+6/7 = 13/7 \ne 1$.
\end{proof}

\begin{remark}
The contradiction is between two clauses of the conjecture itself and does not
require the exact computation of $M(5,2)$: whatever the mean is, the conjecture
demands two different values for it. Theorem~\ref{thm:mean-one} identifies the
correct value, namely $1$.
\end{remark}

\section{The variance clause is false}

Let $N = N(a,b)$ for uniform $(a,b) \in \mathbb{F}_q^2$. We compute the exact
second moment.

\begin{lemma}\label{lem:second-moment}
For every prime power $q$ and every $d$,
\[
\mathbb{E}[N^2] \;=\; 2 - \frac{1}{q}.
\]
\end{lemma}

\begin{proof}
For each unordered pair $\{x,y\}$ of \emph{distinct} points of $\mathbb{F}_q$, we
count the pairs $(a,b)$ for which both $x$ and $y$ are roots:
\[
x^d + ax + b = 0, \qquad y^d + ay + b = 0.
\]
Subtracting the second equation from the first,
\[
(x^d - y^d) + a(x-y) = 0.
\]
Since $x \ne y$, the difference $x-y$ is nonzero and hence invertible in the field
$\mathbb{F}_q$, so
\[
a = -\frac{x^d - y^d}{x-y}
\]
is uniquely determined; substituting back gives $b = -x^d - ax$, also unique.
Thus each unordered pair of distinct points is a pair of common roots for exactly
one $(a,b)$, and there are $\binom{q}{2} = q(q-1)/2$ such pairs.

Now sum $N(a,b)^2$ over all $(a,b)$: this counts ordered pairs $(x,y)$ of roots of
the same trinomial. Splitting into the diagonal $x=y$ and the off-diagonal
$x \ne y$,
\[
\sum_{a,b} N(a,b)^2
= \sum_{a,b} N(a,b) + 2 \cdot \#\{(a,b,\{x,y\}) : x \ne y,\ x,y \text{ roots}\}.
\]
By Theorem~\ref{thm:mean-one}, $\sum_{a,b} N(a,b) = q^2 M(q,d) = q^2$. The
off-diagonal unordered pairs number $\binom{q}{2}$, by the previous paragraph.
Hence
\[
\sum_{a,b} N(a,b)^2 = q^2 + 2\binom{q}{2}
= q^2 + q(q-1) = q(2q-1),
\]
and dividing by $q^2$,
\[
\mathbb{E}[N^2] = \frac{q(2q-1)}{q^2} = \frac{2q-1}{q} = 2 - \frac{1}{q}. \qedhere
\]
\end{proof}

\begin{corollary}\label{cor:variance}
For every prime power $q$,
\[
\operatorname{Var}(N) = \mathbb{E}[N^2] - (\mathbb{E}[N])^2 = 1 - \frac{1}{q}.
\]
In particular the conjectured variance $2 - 1/q$ is false for every $q>1$.
\end{corollary}

\begin{proof}
By Theorem~\ref{thm:mean-one}, $\mathbb{E}[N] = M(q,d) = 1$. Hence, by
Lemma~\ref{lem:second-moment},
\[
\operatorname{Var}(N) = \left(2-\frac1q\right) - 1^2 = 1 - \frac{1}{q}.
\]
For $q>1$ we have $1 - 1/q \ne 2 - 1/q$. Explicitly, $q=3,5,11$ give variances
$2/3$, $4/5$, $10/11$, whereas the conjecture claims $5/3$, $9/5$, $21/11$.
\end{proof}

\begin{remark}
The conjectured variance $2-1/q$ is exactly the second moment
$\mathbb{E}[N^2]$ of Lemma~\ref{lem:second-moment}. Likewise, when $g=1$ the
conjectured ``main term'' $1+(q-1)/q = 2-1/q$ is also exactly this second moment.
So the conjecture has substituted the second moment both for the mean (up to the
$g$-dependent factor) and for the variance.
\end{remark}

\section{The ``main term'' objection, and why it fails}

A charitable reader might propose that $1 + (q-1)/q^{\gcd(d-1,q-1)}$ is asserted
only ``as the main term'', i.e.\ as an asymptotic leading term rather than an exact
finite value, so that Theorem~\ref{thm:mean-one} does not refute it. This reading
does not save the conjecture.

\begin{enumerate}
\item \textbf{The statement is about finite $q$, and the correction is not
negligible.} The definition fixes a single finite field $\mathbb{F}_q$, and the
formula is asserted for that $q$. Even under a leading-term reading, the leading
term of a quantity that is identically $1$ must be $1$. The proposed correction
$(q-1)/q^{g}$ is not a lower-order error: when $g=1$ it equals $(q-1)/q > 1/2$ for
all $q \ge 2$, which is $O(1)$ and cannot be absorbed into lower-order terms. An
asymptotic claim with leading constant $1 + \Omega(1)$ is simply false for a
quantity equal to $1$.
\item \textbf{``Main term'' cannot reconcile two contradictory clauses.} The
conjecture asserts both the value $1+(q-1)/q^{g}$ and, when $g=1$, that the mean
``is exactly $1$''. At $q=5$, $d=2$ this reads $9/5$ versus $1$. No interpretation
of ``main term'' makes one number simultaneously equal to $9/5$ and to $1$.
\item \textbf{The variance clause is separately false.} Independent of any reading
of ``main term'', the assertion ``the variance is $2-1/q$'' is plain arithmetic,
and Corollary~\ref{cor:variance} shows the variance is $1-1/q \ne 2-1/q$ for every
$q>1$.
\end{enumerate}

Therefore all three components fail: the exact mean is $1$, not
$1+(q-1)/q^{g}$; for $g=1$ the value $1+(q-1)/q$ is the second moment, not the
mean; and the variance is $1-1/q$, not $2-1/q$.

\section*{Summary}

For every prime power $q$ and every $d$: $M(q,d) = 1$;
$\mathbb{E}[N^2] = 2 - 1/q$; $\operatorname{Var}(N) = 1 - 1/q$. The conjecture's
main term, exact-mean clause, and variance clause are respectively wrong,
internally contradictory, and wrong. Hence conjecture \texttt{00000001097} is
\textbf{FALSE}.

\end{document}
